\SourceLeaf{445}%
\[
(\xi\xi'+\eta\eta'+\zeta\zeta')\frac{\partial t}{\partial\theta}=\frac{\partial x}{\partial s}.\frac{\partial x}{\partial\theta}+\frac{\partial y}{\partial s}.\frac{\partial y}{\partial\theta}+\frac{\partial z}{\partial s}.\frac{\partial z}{\partial\theta},
\]
wobei wir diese Größe mit $u$ bezeichnen wollen. Sie ist also ihrerseits eine Funktion von $\theta$ und $s$.

Differenziert man nach $s$, so erhält man
\[
\begin{aligned}
\frac{\partial u}{\partial s}={}&\frac{\partial\partial x}{\partial s^{2}}.\frac{\partial x}{\partial\theta}+\frac{\partial\partial y}{\partial s^{2}}.\frac{\partial y}{\partial\theta}+\frac{\partial\partial z}{\partial s^{2}}.\frac{\partial z}{\partial\theta}\\
&+\tfrac12\frac{\partial\left\{\left(\frac{\partial x}{\partial s}\right)^{2}+\left(\frac{\partial y}{\partial s}\right)^{2}+\left(\frac{\partial z}{\partial s}\right)^{2}\right\}}{\partial\theta}\\
={}&\frac{\partial\partial x}{\partial s^{2}}.\frac{\partial x}{\partial\theta}+\frac{\partial\partial y}{\partial s^{2}}.\frac{\partial y}{\partial\theta}+\frac{\partial\partial z}{\partial s^{2}}.\frac{\partial z}{\partial\theta},
\end{aligned}
\]
weil
\[
\left(\frac{\partial x}{\partial s}\right)^{2}+\left(\frac{\partial y}{\partial s}\right)^{2}+\left(\frac{\partial z}{\partial s}\right)^{2}=1,
\]
also sein Differential $=0$ ist.

Da aber alle zu konstantem $\theta$ [gehörigen] Punkte auf einer kürzesten Linie liegen, gilt, wenn $L$ den Zenit des Punktes bezeichnet, dem $s,\theta$ entsprechen, und die Koordinaten von $L$ mit $X,Y,Z$ bezeichnet werden:
\[
\frac{\partial\partial x}{\partial s^{2}}=\frac{X}{p},\qquad\frac{\partial\partial y}{\partial s^{2}}=\frac{Y}{p},\qquad\frac{\partial\partial z}{\partial s^{2}}=\frac{Z}{p},
\]
wobei $p$ den Krümmungsradius bezeichnet. Also gilt
\[
p.\frac{\partial u}{\partial s}=X\frac{\partial x}{\partial\theta}+Y\frac{\partial y}{\partial\theta}+Z\frac{\partial z}{\partial\theta}=\frac{\partial t}{\partial\theta}(X\xi'+Y\eta'+Z\zeta').
\]
Nun ist aber
\[
X\xi'+Y\eta'+Z\zeta'=\cos L\lambda'=0,
\]
weil $\lambda'$ offenbar auf dem Großkreis liegt, dessen Pol $L$ ist. Daher gilt
\[
\frac{\partial u}{\partial s}=0
\]
oder $u$ ist von $s$ unabhängig und daher nur eine Funktion von $\theta$. Für $s=0$ gilt aber offenbar $t=0$, $\frac{\partial t}{\partial\theta}=0$ und somit $u=0$. Daraus schließen wir, dass allgemein $u=0$ gilt oder
\[
\cos\lambda\lambda'=0.
\]
Daraus folgt der schöne Lehrsatz:
